\chapter[Evaluating a real-company simulator]{Evaluating a real-company simulator: what evidence would support its use}
\label{sim:evaluation}

The exact cases in this part establish useful semantics, but a real-company simulator needs a broader evaluation. It must represent disclosures, preserve important uncertainty, generate coherent events and provide decision-relevant results. These objectives cannot be reduced to one extraction score or one forecast error. This chapter proposes an evaluation design and explains what each result would and would not establish. No empirical evaluation is claimed to have been performed here.

\section{Define the unit of evaluation}
A unit can be a source assertion, a contract, a reporting period, a company history family, a scenario query or a policy decision. Accuracy at one unit does not automatically transfer to another. Correctly extracting a debt total does not establish that maturity and option terms are represented.

A company-period case should identify its document package, entity perimeter, accounting basis, evidence cutoff and materiality scope. A scenario case adds shocks, policies, information and requested outputs. These definitions prevent evaluation from quietly changing the target after seeing results.

The unit should also preserve unresolved content. If a system refuses an unsupported instrument, that can be correct behavior within a bounded profile. The evaluation should distinguish honest refusal from silent omission, while still recording the coverage limitation for the full mission.

\section{Source-to-model traceability}
For each material modeled assertion, a reviewer should be able to reach the source passage and interpretation. Conversely, material source passages should have a disposition: represented, irrelevant to the declared query, unresolved or outside the executable profile.

This bidirectional trace catches different failures. A model assertion without a source may be an assumption or hallucination. A source passage without a disposition may be an omission. A correct numeric total can coexist with a missed covenant note that changes the stress result.

The evaluation should sample across difficult formats and concepts, including scanned notes, tables with overlapping dimensions, incorporated documents and revisions. A convenient set of clean tables is insufficient evidence for whole-filing coverage. The selection procedure and omissions should be recorded.

\section{Annotation and disagreement}
Human annotations of financial meaning can disagree. A term may be ambiguous, a source may conflict, or the accounting treatment may depend on facts not disclosed. The reference set should permit unresolved or alternative annotations rather than force artificial unanimity.

Adjudication should record the competing interpretations and the evidence that resolves them. A majority vote is not necessarily an authority on a technical accounting or contractual question. For some cases the correct reference is a conditional family.

Ordinary comprehension review and expert technical review serve different purposes. A reader can identify unclear exposition without certifying accounting treatment. An accountant can assess a specific treatment without validating a probabilistic forecasting model. Evaluation should not collapse these roles into a generic human-approved label.

\section{Historical reconstruction checks}
A reconstruction should reproduce observed statements within admitted precision and context. It should also satisfy cross-statement bridges and contract constraints. These are necessary checks, but a model explicitly fitted to the reports can pass them while selecting the wrong hidden history.

Use withheld disclosures where possible. For example, reconstruct from primary statements and selected notes, then test whether the retained family includes a contract detail revealed in a withheld exhibit. The procedure must avoid using that detail indirectly through another source.

A family that excludes the revealed truth under correct source interpretation has failed coverage. A family that includes everything can pass inclusion trivially but be uninformative. Evaluation therefore needs both coverage and useful restriction, interpreted relative to the available evidence.

\section{Synthetic truth and real uncertainty}
Synthetic companies provide known complete histories and controlled disclosure maps. They allow tests of whether the system preserves all intended alternatives and rejects impossible ones. They are valuable for semantic coverage and adversarial construction.

Their limitation is that the generator defines the world. A system matching that generator may fail on mechanisms the generator omits. Synthetic results should not be presented as real-company empirical validity. The generator's assumptions and diversity need review.

Real disclosures supply realistic ambiguity but often lack a fully observed ground truth. Evaluation can then assess traceability, consistency, withheld facts, conditional predictions and decision consequences. The two evidence sources are complementary rather than interchangeable.

\section{Forecast origins and leakage}
A forecast should be reconstructed using information available at its origin. Later restatements, realized prices and revised macro data can leak into features or labels. The evidence store needs dated versions to prevent this.

A rolling-origin evaluation can produce forecasts at several historical dates and compare them with subsequent observations. The model's parameter estimation and tuning must also respect the cutoff. Training on the full sample and then calling earlier predictions out of sample is incorrect.

Company events such as acquisitions or changes of reporting basis can alter comparability. The evaluation should bridge those changes or define a target that remains meaningful. A forecast error caused by a perimeter mismatch is not the same as an economic prediction error.

\section{Conditional scenario evaluation}
A deterministic stress scenario may never occur exactly. Its mechanical correctness can still be tested through contract formulas, accounting invariants and controlled cases. Empirical validation of conditional responses requires suitable observed variation or another identification strategy.

A model predicting cash under a stipulated 200-basis-point shock cannot be validated by comparing it with actual cash after a period containing many other shocks without accounting for those differences. Such a comparison may be descriptive, but it does not isolate the rate response.

The evaluation should distinguish unconditional forecasting, conditional stress mechanics and causal intervention claims. Each has a different comparator and evidence requirement. Calling all three simulation accuracy obscures the target.

\section{Probabilistic forecasts}
If the system supplies a predictive distribution, evaluate both calibration and informativeness with methods appropriate to the outcome. A forecast assigning broad mass everywhere may be calibrated in a weak sense while offering little value. A sharp forecast can be badly miscalibrated.

Proper scoring rules, discussed in \citet{gneiting2007proper}, provide a principled way to reward honest predictive distributions under their conditions. The chosen score should match the outcome type and decision context. A score on annual earnings does not automatically validate default probability or liquidity tails.

Probabilities conditional on one reconstructed history should be distinguished from probabilities averaged across histories. If history weights are subjective or ambiguous, report sensitivity. A good score under one prior does not establish that the prior is uniquely justified.

\section{Set-valued forecasts}
An unweighted family produces a set of possible outcomes. Evaluation can ask whether later observed outcomes fall within it and how wide or structured the set is. A rectangular interval display may obscure joint restrictions, so multivariate coverage should respect the actual feasible set.

If observed cash and interest each lie within their marginal intervals but violate the modeled relation $C+I=95$, the joint model has failed for that observation. Marginal coverage alone would miss it. Conversely, a wide marginal interval may coexist with a precise relation that is useful for a decision.

A set can be evaluated against a declared decision task: does it correctly identify policies that are robustly feasible or expose the history driving a failure? This should supplement, not replace, source and mathematical checks. A useful decision result can still rest on an incomplete family.

\section{Policy evaluation}
A policy should be evaluated under the information and action constraints it would actually face. A perfect-foresight benchmark can provide an upper bound in some settings, but it is not an implementable comparator. A no-action baseline may be useful but can be too weak to establish practical value.

A fair comparison includes a credible existing policy or classical method, with equivalent access to data and tuning. If the proposed method receives a richer contract representation, separate the benefit of representation from the benefit of its optimizer. Otherwise a solver improvement claim may actually be a data improvement.

Policy value estimates need uncertainty analysis and attention to rare failures. A higher mean payout can come with unacceptable liquidity risk. The primary criterion and vetoes should be declared before selecting a winner. Descriptive advantages from a few simulations should not be promoted to a default policy.

\section{Ablations that answer real questions}
An ablation can remove contract detail, actor information, persistent memory or joint constraints to test their contribution. The comparison should preserve other conditions. Removing several features at once can show that the bundle matters but not which mechanism explains the change.

For example, compare a full debt schedule with an aggregate debt balance while holding pricing and policy fixed. Then test the effect of adding endogenous refinancing. This sequence separates contract-state sufficiency from behavioral response.

Expected failures can be informative. If the aggregate model misses a maturity wall, that supports the need for timing in that query. It does not prove that every query requires instrument-level detail. The conclusion should match the demonstrated mechanism and scope.

\section{Stress coverage}
A validation set should include ordinary operations, near-boundary cases and severe but admitted stresses. Important dimensions include rates, currency, customer concentration, maturity, collateral, offer capacity, settlement delay and information timing.

A Cartesian product of every dimension may be infeasible or economically incoherent. Use the joint model to identify feasible combinations and targeted adversarial cases. Cover interactions known to matter, such as collateral loss reducing funding exactly when cash falls.

Coverage should record mechanisms tested and regions untested. A large path count is not a substitute. Ten thousand smooth scenarios can leave a single conversion or cure-period boundary unchecked.

\section{Numerical reliability}
For approximate methods, validation should compare against exact or highly controlled small references, assess convergence or refinement, and inspect downstream decisions. Solver tolerances should be translated into economic quantities where possible.

A residual below a threshold can still be unacceptable if it changes default classification. A slightly larger residual in a large nonbinding balance may be immaterial. The evidence contract should state which diagnostics veto use and which explain remaining error.

Repeat stochastic comparisons with appropriate seeds and uncertainty estimates. If candidates remain statistically indistinguishable after validity screens, report that. A method can be viable without being proven superior. Default promotion requires stronger evidence than an optional experimental path.

\section{Failure accounting}
Evaluation must count unsupported, refused and numerically failed cases. Excluding them from average error can make a fragile method look accurate on the easy subset. Report coverage and conditional performance together.

A refusal can be correct if the request exceeds the admitted profile, but excessive refusal limits usefulness. A numerical failure is different from a correctly identified economic shortfall. A result table should preserve these categories rather than assign all failures the same placeholder number.

If an approximation leaves probability mass unresolved, include that mass in risk interpretation. If the model is unweighted, report unresolved regions and the bounds they permit. Renormalizing completed cases is a different conditional target and should be labeled as such.

\section{Reproducibility and independent checks}
A reproducible run records source snapshots, model versions, assumptions, methods, seeds and environment. It also preserves exact commands and outputs. Reproduction verifies that the same computation can be obtained; it does not independently validate the model.

Independent checks should use different routes where feasible: algebra versus code, event-level versus aggregate reconciliation, or one solver versus a certified reference. Shared data and code can share errors. A second model reading the same incorrect assumption is not independent economic evidence.

Rendered document review is also distinct. It can reveal missing symbols, broken references and unreadable derivations. A clean build cannot establish that a reader understands the argument. The manuscript should remain provisional until actual human comprehension evidence is available.

\section{A proposed staged evaluation}
First establish exact semantics with the synthetic cases in this part. Then evaluate source-to-model coverage on a bounded real disclosure package, preserving unresolved content. Next test historical reconstruction against withheld facts and carefully defined reporting bridges.

Only after those stages should a forecast or policy study claim economic performance. Its plan must identify target companies, time cutoffs, baselines, tuning protocol, primary criterion and uncertainty analysis. A failure at an earlier stage can invalidate interpretation of the later results.

The stages are dependencies, not excuses to stop independent work. Source research, documentation and bounded implementation can continue while empirical data or human review remain pending. The full mission remains open until its specified evidence exists.

\section{What a defensible conclusion sounds like}
A strong narrow conclusion might state that a method reproduces all exact cash and ownership transitions in a declared profile and preserves the tested information restrictions. Another might state that a predictive distribution outperforms a named baseline under a predeclared score with uncertainty evidence on a stated sample.

Neither conclusion should be expanded into universal accounting compliance or company forecasting ability. A failed candidate can be rejected without rejecting the entire uncertainty-preserving approach. A successful small model can justify the next experiment without becoming a production default.

The evaluation program should make those distinctions easy to maintain. That is the practical counterpart of the manuscript's central idea: preserve what is known, preserve what remains possible, and make every stronger conclusion earn its evidence.
